%%=============================================================================
%% Bijlage: implementatiecode
%%=============================================================================

\chapter{Implementatiecode}%
\label{ch:bijlage-implementatie}

Deze bijlage bevat de volledige codefragmenten die in Hoofdstuk~\ref{ch:implementatie} cherry-picked werden. De volgorde is dezelfde als die van het hoofdstuk, zodat de lezer per onderwerp eenvoudig kan terugbladeren.

\section{Orb agent: YAML-configuraties}
\label{sec:bijlage-orb-yaml}

\subsection{Device discovery}

\begin{listing}[H]
  \caption[Orb agent voor device discovery]{Volledige Orb-agentconfiguratie voor device discovery via Diode (\S\ref{sec:architectuur})}
  \label{lst:orb-device-discovery-config}
  \begin{minted}[fontsize=\small]{yaml}
orb:
  log_level: DEBUG
  config_manager:
    active: local

  backends:
    common:
      diode:
        target: grpc://localhost:8080/diode
        client_id: ${DIODE_CLIENT_ID}
        client_secret: ${DIODE_CLIENT_SECRET}
        agent_name: discovery_agent_01
        dry_run: true
        dry_run_output_dir: /opt/orb
    device_discovery:
      host: 0.0.0.0
      port: 8072
      log_level: DEBUG
      log_format: TEXT

  policies:
    device_discovery:
      oostkamp_router:
        config:
          schedule: "*/30 * * * *"
          options:
            port_scan_ports: [8722]
            port_scan_timeout: 1.0
            debug: true
          defaults:
            site: "Oostkamp Ridefort"
            role: "Router"
            tenant: "Oostkamp"
        scope:
          - hostname: "10.23.4.225"
            driver: "ros"
            username: "${DEVICE_USERNAME}"
            password: "${DEVICE_PASSWORD}"
            optional_args:
              port: 8722
              use_ssl: false
              debug: true
  \end{minted}
\end{listing}

\subsection{SNMP discovery}

\begin{listing}[H]
  \caption[Orb agent voor SNMP discovery]{Volledige Orb-agentconfiguratie voor SNMP discovery (\S\ref{sec:architectuur})}
  \label{lst:orb-snmp-discovery-config}
  \begin{minted}[fontsize=\small]{yaml}
orb:
  config_manager:
    active: local

  backends:
    common:
      diode:
        target: grpc://127.0.0.1:8080/diode
        client_id: ${DIODE_CLIENT_ID}
        client_secret: ${DIODE_CLIENT_SECRET}
        agent_name: discovery_agent_01

  policies:
    snmp_discovery:
      oostkamp_ridefort:
        config:
          schedule: "*/30 * * * *"
          timeout: 300
          snmp_timeout: 5
          retries: 3
          lookup_extensions_dir: "/opt/orb/snmp-extensions/lookup_extensions"
          defaults:
            site: "Oostkamp Ridefort"
            role: "network"
            rack: "Oostkamp Ridefort"
            tags: ["snmp-discovery", "orb"]
            device:
              tenant: "Oostkamp"
              rack: "oostkamp_ridefort"
              description: "SNMP discovered device"

        scope:
          targets:
            - host: "10.23.4.225"
              override_defaults:
                role: "Router"

            - host: "10.23.4.226"
              override_defaults:
                role: "Switch"

            - host: "10.23.4.231-236"
              override_defaults:
                role: "Access Point"

          authentication:
            protocol_version: "SNMPv2c"
            community: "public"
  \end{minted}
\end{listing}

In bovenstaande YAML- en shellfragmenten staat de site-naam in NetBox als \texttt{Oostkamp Ridefort}; in de hoofdtekst wordt deze site kortweg \emph{Ridefort} genoemd. Credentials zoals \texttt{DEVICE\_USERNAME}, \texttt{DEVICE\_PASSWORD} en de SNMP-community \texttt{public} dienen in deze listings als placeholders en worden in productie via een secret manager respectievelijk een niet-publieke community vervangen.

\section{Orb agent: shellscripts voor uitrol}
\label{sec:bijlage-orb-scripts}

\subsection{Script voor SNMP discovery}

\begin{listing}[H]
  \caption[Shellscript voor SNMP-discovery uitrol]{Volledig shellscript voor het genereren van \texttt{agent.yaml} voor SNMP discovery (\S\ref{sec:schaalbaarheid})}
  \label{lst:orb-snmp-script}
  \begin{minted}[fontsize=\small]{bash}
#!/usr/bin/env bash
set -e

SCRIPT_DIR="$(cd "$(dirname "${BASH_SOURCE[0]}")" && pwd)"
cd "$SCRIPT_DIR"

mkdir -p "$SCRIPT_DIR/snmp-extensions"

CREDS_FILE="/opt/diode/oauth2/client/client-credentials.json"
ALT_CREDS_FILE="/vagrant/diode/oauth2/client/client-credentials.json"

if [ ! -f "$CREDS_FILE" ]; then
  if [ -f "$ALT_CREDS_FILE" ]; then
    CREDS_FILE="$ALT_CREDS_FILE"
  else
    echo "Fout: $CREDS_FILE niet gevonden."
    exit 1
  fi
fi

if ! command -v envsubst >/dev/null 2>&1; then
  echo "envsubst niet gevonden. Installeer met: sudo apt-get install -y gettext-base"
  exit 1
fi

export DIODE_CLIENT_ID="diode-ingest"
export DIODE_CLIENT_SECRET=$(sudo jq -r '.[] | select(.client_id == "diode-ingest") | .client_secret' "$CREDS_FILE")

if [ -z "$DIODE_CLIENT_SECRET" ]; then
  echo "Fout: kon diode-ingest secret niet ophalen."
  exit 1
fi

export DEVICE_USERNAME="user"
export DEVICE_PASSWORD="pass"

envsubst '${DIODE_CLIENT_ID} ${DIODE_CLIENT_SECRET} ${DEVICE_USERNAME} ${DEVICE_PASSWORD}' < agent.yaml.snmp > agent.yaml
echo "agent.yaml gegenereerd met credentials."

exec sudo docker run --rm --network host \
  -v "$SCRIPT_DIR:/opt/orb" \
  -e PIP_BREAK_SYSTEM_PACKAGES=1 \
  netboxlabs/orb-agent:latest run -c /opt/orb/agent.yaml
  \end{minted}
\end{listing}

\subsection{Script voor router/device discovery}

\begin{listing}[H]
  \caption[Shellscript voor router-discovery uitrol]{Volledig shellscript voor het genereren van \texttt{agent.yaml} voor router/device discovery (\S\ref{sec:schaalbaarheid})}
  \label{lst:orb-router-script}
  \begin{minted}[fontsize=\small]{bash}
#!/usr/bin/env bash
set -e

SCRIPT_DIR="$(cd "$(dirname "${BASH_SOURCE[0]}")" && pwd)"
cd "$SCRIPT_DIR"

mkdir -p "$SCRIPT_DIR/snmp-extensions"

CREDS_FILE="/opt/diode/oauth2/client/client-credentials.json"
ALT_CREDS_FILE="/vagrant/diode/oauth2/client/client-credentials.json"

if [ ! -f "$CREDS_FILE" ]; then
  if [ -f "$ALT_CREDS_FILE" ]; then
    CREDS_FILE="$ALT_CREDS_FILE"
  else
    echo "Fout: $CREDS_FILE niet gevonden."
    exit 1
  fi
fi

if ! command -v envsubst >/dev/null 2>&1; then
  echo "envsubst niet gevonden. Installeer met: sudo apt-get install -y gettext-base"
  exit 1
fi

export DIODE_CLIENT_ID="diode-ingest"
export DIODE_CLIENT_SECRET=$(sudo jq -r '.[] | select(.client_id == "diode-ingest") | .client_secret' "$CREDS_FILE")

if [ -z "$DIODE_CLIENT_SECRET" ]; then
  echo "Fout: kon diode-ingest secret niet ophalen."
  exit 1
fi

export DEVICE_USERNAME="user"
export DEVICE_PASSWORD="pass"

envsubst '${DIODE_CLIENT_ID} ${DIODE_CLIENT_SECRET} ${DEVICE_USERNAME} ${DEVICE_PASSWORD}' < agent.yaml.router > agent.yaml
echo "agent.yaml gegenereerd met credentials."

exec sudo docker run --rm --network host \
  -v "$SCRIPT_DIR:/opt/orb" \
  -e INSTALL_DRIVERS_PATH=/opt/orb/drivers-router.txt \
  -e PIP_BREAK_SYSTEM_PACKAGES=1 \
  netboxlabs/orb-agent:2.5.0 run -c /opt/orb/agent.yaml
  \end{minted}
\end{listing}

\section{NetworkTool: GUI-architectuur}
\label{sec:bijlage-networktool-gui}

\begin{listing}[H]
  \caption[NetworkTool: module-imports en tabstructuur]{NetworkTool: module-imports en tabstructuur (\S\ref{sec:architectuur})}
  \label{lst:networktool-gui-tabs-full}
  \begin{minted}[fontsize=\small]{python}
from config import AGENT_MAP, AGENT_DIR, MAPPING_FILE
from utils import load_snmp_mapping, parse_ip_input, sweep_ip, guess_role_and_type
from network_core import fetch_live_data
from netbox_api import (
    get_nb_sites, get_nb_devices, get_nb_roles,
    get_nb_device_types, get_nb_tenants,
    evaluate_and_sync_netbox, sync_snmp_devices
)

tab_agent, tab_scanner, tab_discovery, tab_cve = st.tabs(
    ["Agent Manager", "Network Scanner", "SNMP Discovery", "CVE Scanner"]
)
  \end{minted}
\end{listing}

\section{NetworkTool: dry-run en live synchronisatie}
\label{sec:bijlage-networktool-sync}

\begin{listing}[H]
  \caption[NetworkTool: dry-run gevolgd door live synchronisatie]{NetworkTool: dry-run-analyse gevolgd door live synchronisatie (\S\ref{sec:schaalbaarheid})}
  \label{lst:networktool-dryrun-live}
  \begin{minted}[fontsize=\small]{python}
success, changes = evaluate_and_sync_netbox(
    live_data_payload, d['ip'], d['site'], d['role'], d['rack'], d['tenant'], d['naam'],
    nb_url, nb_token, selected_modules, d['vendor'], d['monitoring_env'], dry_run=True
)

success, messages = evaluate_and_sync_netbox(
    live_data_payload, d['ip'], d['site'], d['role'], d['rack'], d['tenant'], d['naam'],
    nb_url, nb_token, st.session_state.selected_modules, d['vendor'], d['monitoring_env'],
    dry_run=False
)
  \end{minted}
\end{listing}

\section{NetworkTool: LLDP-verwerking en kabel-drift}
\label{sec:bijlage-networktool-lldp}

\subsection{LLDP ophalen via NAPALM-fallback}

\begin{listing}[H]
  \caption[NetworkTool: LLDP ophalen en formatteren]{NetworkTool: LLDP ophalen en formatteren met fallback (\S\ref{sec:issues-valkuilen})}
  \label{lst:networktool-lldp-fetch}
  \begin{minted}[fontsize=\small]{python}
try:
    details = device.get_lldp_neighbors_detail()
    if details:
        fmt_lldp = {}
        for iface, neighbors in details.items():
            fmt_lldp[iface] = [
                {'hostname': n.get('remote_system_name', ''), 'port': n.get('remote_port', '')}
                for n in neighbors
            ]
        data["lldp_neighbors"] = fmt_lldp
    else:
        data["lldp_neighbors"] = device.get_lldp_neighbors()
except Exception:
    try:
        data["lldp_neighbors"] = device.get_lldp_neighbors()
    except Exception:
        pass
  \end{minted}
\end{listing}

\subsection{Kabel-drift detecteren en corrigeren}

\begin{listing}[H]
  \caption[NetworkTool: kabel-drift detecteren en corrigeren]{NetworkTool: kabel-drift detecteren en corrigeren in NetBox (\S\ref{sec:issues-valkuilen})}
  \label{lst:networktool-cable-drift}
  \begin{minted}[fontsize=\small]{python}
if cable_a:
    endpoints = getattr(nb_intf_a, 'connected_endpoints', [])
    if endpoints and len(endpoints) > 0:
        ep = endpoints[0]
        connected_dev_name = getattr(getattr(ep, 'device', None), 'name', 'Onbekend')
        connected_port_name = getattr(ep, 'name', 'Onbekend')
        if str(connected_dev_name).upper() != str(remote_host).upper() or \
           str(connected_port_name).upper() != str(nb_intf_b.name).upper():
            drift_detected = True

if not kabel_bestaat_al_goed:
    nb.dcim.cables.create({
        "a_terminations": [{"object_type": "dcim.interface", "object_id": nb_intf_a.id}],
        "b_terminations": [{"object_type": "dcim.interface", "object_id": nb_intf_b.id}],
        "status": "connected"
    })
  \end{minted}
\end{listing}

\section{NetworkTool: RUCKUS vSZ-koppeling}
\label{sec:bijlage-vsz}

\begin{listing}[H]
  \caption[RUCKUS vSZ: inventarisdict richting NetBox]{RUCKUS vSZ: inventarisdict richting NetBox (\S\ref{sec:issues-valkuilen})}
  \label{lst:ruckus-vsz-netbox-inventory}
  \begin{minted}[fontsize=\small]{python}
data = {
    "name": ap_name,
    "device_type": dev_type.id,
    "role": role.id,
    "site": site.id,
    "serial": ap_serial,
    "status": "active",
    "tenant": tenant.id if tenant else None,
    "asset_tag": ap_mac
}

if ap_vlan:
    data["description"] = f"Management VLAN: {ap_vlan}"

# Vervolgens wordt het device aangemaakt of geüpdatet met data,
# inclusief IP-adresverwerking via deze vSZ-stap.
  \end{minted}
\end{listing}
